\begin{table}[htbp]
\centering
\caption{\textbf{Extended and sensitivity models (KNHANES, adults 40--69).}
Perceived-stress odds ratio (95\% CI) across prespecified specifications,
isolating the robustness of the primary exposure. Rows add the mediator
(depressed mood) or audiometric hearing loss to the minimal-sufficient set
(perceived stress, occupational noise, age, sex), and vary the hearing-loss
outcome definition (better-ear, worse-ear, high-frequency). All are design-based
associations, not causal effects; mediator-adjusted rows are interpreted as
mediator- (not confounder-) adjusted. Generated by
\texttt{code/src/knhanes\_analysis.py}. Cells marked \emph{[pending data run]}
populate automatically when the pipeline is executed on the KDCA-approved
microdata with \texttt{-{}-latex-extended}. $^{*}p<0.05$, $^{**}p<0.01$,
$^{***}p<0.001$.}
\label{tab:results_knhanes_extended}
\small
\begin{tabular}{lcc}
\toprule
Model specification & Perceived stress, OR (95\% CI) & $N$ \\
\midrule
Tinnitus --- minimal-sufficient (primary) & 1.42 (1.26--1.60)*** & 9504 \\
\quad + depressed mood (mediator-adjusted) & \textit{[pending data run]} & -- \\
\quad + audiometric hearing loss & \textit{[pending data run]} & -- \\
Hearing loss, better ear --- minimal-sufficient & 1.18 (0.99--1.41) & 9176 \\
\quad + depressed mood (mediator-adjusted) & \textit{[pending data run]} & -- \\
Hearing loss, worse ear --- minimal-sufficient & \textit{[pending data run]} & -- \\
Hearing loss, high-freq (3/4/6\,kHz) --- minimal-sufficient & \textit{[pending data run]} & -- \\
Bothersome tinnitus --- minimal-sufficient & \textit{[pending data run]} & -- \\
Non-bothersome tinnitus vs none (reverse-causation sensitivity) & \textit{[pending data run]} & -- \\
\bottomrule
\end{tabular}
\end{table}
